\documentclass[addpoints]{exam}
% \documentclass[addpoints,12pt]{exam}

\usepackage[slovene]{babel}
\usepackage{amsfonts}
\usepackage{amssymb}
\usepackage{amsmath}
\usepackage{float}
\usepackage{graphicx}
\usepackage{enumitem}
\usepackage{color}
\usepackage{eurosym}

\usepackage{pgfpages}
% \usepackage{tikz}
\usepackage{wrapfig}
\usepackage{pgfkeys}
\usepackage{pgfplots}
\usepackage{xcolor}
\usepackage{tkz-euclide}
\usepackage{xfp}




%Komentar naslednje vrstice glede na verzijo testa -- rešitve ali prostor za odgovore:
% \printanswers

% \linespread{2}


\pointsinrightmargin
\marginbonuspointname{}


\renewcommand{\solutiontitle}{\noindent\textbf{Rešitve in točkovanje:}\par\noindent}

\colorgrids
\definecolor{GridColor}{gray}{0.7}
\setlength{\gridsize}{7mm}

\extrawidth{5mm}
\setlength{\rightpointsmargin}{1.5cm}

\begin{document}

\runningfooter{}{}{}

\begin{center}
    \fbox{\fbox{\parbox{15cm}{\centering
    Popravljanje in pridobivanje ocen -- konferenca \\ 2. letnik -- test G \\ 23. januar 2026 \\ (Geometrija v ravnini; Kotne funkcije; Vektorji)}}}
\end{center}

\vspace{0.1in}
\makebox[\textwidth]{Ime in priimek, oddelek:\enspace\hrulefill}


\begin{center}
    \hqword{Naloga}
    \hpword{Možne točke}
    % \chbpword{Dodatne točke}
    \hsword{Dosežene točke}
    \htword{Skupaj}  
    % \combinedpointtable[h][questions]
    \gradetable[h][questions]

\end{center}

\vspace{0.1in}


\begin{questions}

    % 1. vprašanje

    \question[5]
        Izračunajte velikosti vseh notranjih in zunanjih kotov trikotnika $\triangle ABC$, 
        če je vsota velikosti dveh zunanjih kotov $\alpha'+\beta'=234^\circ$, razlika velikosti dveh notranjih kotov pa $\alpha-\beta=28^\circ$.

    \begin{solution}[10cm]
        $ \dfrac{n(n-3)}{2}=44 \Rightarrow n^2-3n-88=0 \Rightarrow (n-11)(n+8)=0 \Rightarrow n=11 $ \\
        $ \varphi=\dfrac{n-2}{n}180^\circ=\dfrac{1620}{11}^\circ=147.\overline{27}^\circ $ \\ \\
        Za zapis in pravilno uporabo formule za število diagonal $1$ točka, za pravilen izračuna $1$ točka. 
        Za pravilen zapis in uporabo formule za izračun notranjega kota pravilnega $n$-kotnika $1$ točka, za pravilen izračun $1$ točka.
    \end{solution}


    % % 2. vprašanje

    % \question[3]
    % Središčni kot je za $64^\circ$ večji od obodnega kota nad istim lokom. 
    % Izračunajte velikosti obeh kotov.

    % \begin{solution}[2cm]
    %     $\varphi+\psi=180^\circ$ \\
    %     $\varphi-\psi=25^\circ 29' 48''$ \\
    %     $\varphi=102^\circ 44' 54''$ in $\psi=77^\circ 15' 6 ''$ \\ \\
    %     Za pravilno postavitev enačb sistema po $1$ točka, reševanje sistema $1$ točka, za pravilen izračun rešitev po $1$ točka.
    % \end{solution}



    \question[3]
        Vsota števila stranic in diagonal $n$-kotnika je $105$. Izračunajte, kateri $n$-kotnik je to.

    \begin{solution}[4cm]
        Uporabimo lahko Pitagorov izrek za dolžini kraka in njegove pravokotne projekcije na osnovnico: $v^2=b^2-\left(\frac{a-c}{2}\right)^2=17^2-8^2=225 \Rightarrow v=15~cm$ \\ \\
        Za pravilen razmislek in izračun posameznih dolžin $2$ točki, pravilna uporaba Pitagorovega izreka $1$ točka, pravilen izračun $1$ točka.
    \end{solution}


    \newpage
    % 3. vprašanje

    \question[6]
    Konstruirajte trapez s podatki $a=5~cm$, $e=f=4.5~cm$ in $\alpha=60^\circ$.

    \begin{solution}[18cm]
        Za zapis konstrukcije s pravilnim izrisom konstrukcije do $6$ točk. Zgolj za izris konstrukcije $1$ točka.
    \end{solution}




    % \newpage
    % % 4. vprašanje

    % \question[5]
    %     S skice preberite ustrezne podatke ter izračunajte velikosti kotov $\alpha$, $\beta$, $\gamma$ in $\delta$.
    %     Pri tem velja, da so premice $p$, $q$ in $r$ vzporedne.
    %     \begin{figure}[H]
    %         \centering
    %         \begin{tikzpicture}
    %             % \clip (0,0) rectangle (14.000000,10.000000);
    %             {\footnotesize

    %             % Drawing line Q
    %             \draw [line width=0.016cm] (1.000000,3.000000) -- (5.800000,3.000000);%

    %             % Drawing line P
    %             \draw [line width=0.016cm] (1.000000,4.000000) -- (5.800000,4.000000);%

    %             % Drawing line R
    %             \draw [line width=0.016cm] (1.000000,1.500000) -- (5.800000,1.500000);%

    %             % Drawing line S
    %             \draw [line width=0.016cm] (2.426509,1.000000) -- (3.516142,4.800000);%

    %             % Drawing line T
    %             \draw [line width=0.016cm] (4.865030,1.000000) -- (1.321473,4.800000);%

    %             % Marking point p
    %             \draw (5.500000,4.000000) node [anchor=south] { $p$ };%

    %             % Marking point q
    %             \draw (5.500000,3.000000) node [anchor=south] { $q$ };%

    %             % Marking point r
    %             \draw (5.500000,1.500000) node [anchor=south] { $r$ };%

    %             % Marking point s
    %             \draw (3.500000,4.700000) node [anchor=west] { $s$ };%

    %             % Marking point t
    %             \draw (1.500000,4.700000) node [anchor=west] { $t$ };%

    %             % Marking point \alpha
    %             \draw (1.800000,4.000000) node [anchor=south east] { $\alpha$ };%

    %             % Marking point \beta
    %             \draw (3.300000,3.000000) node [anchor=north west] { $\beta$ };%

    %             % Marking point \gamma
    %             \draw (3.000000,3.000000) node [anchor=north east] { $\gamma$ };%

    %             % Marking point \delta
    %             \draw (4.450000,1.500000) node [anchor=south] { $\delta$ };%

    %             % Marking point {x-15^\circ}
    %             \draw (2.000000,3.000000) node [anchor=south] { ${x-15^\circ}$ };%

    %             % Marking point {2x-18^\circ}
    %             \draw (4.000000,4.000000) node [anchor=north] { ${2x-18^\circ}$ };%

    %             % Marking point {x+12^\circ}
    %             \draw (3.250000,1.500000) node [anchor=south] { ${x+12^\circ}$ };%
    %             }
    %         \end{tikzpicture}

    %     \end{figure}

    % \begin{solution}[4cm]
    %     $x=65^\circ$ -- obodni kot nad lokom $AD$ \\
    %     $z=180^\circ-90^\circ-65^\circ=25^circ$ -- kot v pravokotnem trikotniku $\triangle ACD$ \\
    %     $y=180^\circ-65^\circ-25^\circ-35^\circ=55^\circ$ -- kot v trikotniku $\triangle ABD$ \\ \\
    %     Za izračun kota $x$ in utemeljitev $1$ točka, za izračun kota $y$ in utemeljitev $1$ točka, upoštevanje Talesovega izreka v polkrogu $1$ točka, izračuna kota $z$ $1$ točka.
    % \end{solution}



    % 5. vprašanje

    % \question[4]
    %     Izračunajte dolžine stranic pravokotnega trikotnika $\triangle ABC$, kjer sta $a$ in $b$ kateti, $c$ pa hipotenuza:
    %     $a=3x-1$, $b=5x$, $c=6x-1$. Določite dolžino pravokotne projekcije katete $a$ na hipotenuzo.


    %     \begin{solution}[9cm]
    %         Z upoštevanjem lastnosti pravokotnega trikotnika je $t_c=\frac{c}{2}=6~cm$. \\
    %         Za pravilen izračun dolžine $1$ točka.
    %     \end{solution}




    % \newpage
    % 6. vprašanje

    \question[2]
        Marija stoji ponoči $5~m$ od vznožja $9~m$ visokega svetilnika. Njena senca je dolga $1~m$.
        Določite Marijino višino.


    \begin{solution}[2cm]
        Uporabimo lahko Pitagorov izrek za dolžini kraka in njegove pravokotne projekcije na osnovnico: $v^2=b^2-\left(\frac{a-c}{2}\right)^2=17^2-8^2=225 \Rightarrow v=15~cm$ \\ \\
        Za pravilen razmislek in izračun posameznih dolžin $2$ točki, pravilna uporaba Pitagorovega izreka $1$ točka, pravilen izračun $1$ točka.
    \end{solution}




\newpage
 % 1. vprašanje

    \question[6]
    Poenostavite izraz in izračunajte njegovo natančno vrednost. 
    $$\displaystyle \left(\left(\tan{x}\cos{x}\right)^{-2}+\cos^{-2}{x}\right)\sin^2{x}\cos^2{x} + \left(\cos{135^\circ}+\sin{225^\circ}\right)^2$$

    \begin{solution}[12cm]
       
    \end{solution}


    % % 2. vprašanje

    % \question[5]
    % Izračunajte natančno vrednost izraza. 
    % $$\displaystyle \dfrac{\left(\cos{135^\circ}+\sin{225^\circ}\right)^2}{\tan{300^\circ}-\tan{120^\circ}-\sin{270^\circ}} $$

    % \begin{solution}[6cm]
       
    % \end{solution}


% \newpage
    % 3. vprašanje

    \question[3]
    V pravilnem šestkotniku $ABCDEF$ točka $M$ tretjini stranico $BC$, točka $N$ razdeli stranico $ED$ v razmerju $|EN|:|ND|=2:3$.
    Z baznima vektorjema $\overrightarrow{a}=\overrightarrow{AB}$ in $\overrightarrow{b}=\overrightarrow{AF}$ izrazite vektor $\overrightarrow{MN}$. 
    % \begin{parts}
    %     \part[3]
    %     $\overrightarrow{MN}$, 

    % \begin{solution}[5cm]

    % \end{solution}
        
    %     \part[2] 
    %     $\overrightarrow{NB}-\overrightarrow{FB}+\overrightarrow{FA}-\overrightarrow{EA}$.
    % \end{parts}

    \begin{solution}[4cm]

    \end{solution}

    

    
    % % 4. vprašanje

    % \question[5]
    % V trikotniku $\triangle ABC$ z oglišči $A(4,5,-6)$, $B(1,7,0)$ in $C(4,0,-3)$ izračunajte oddaljenost težišča trikotnika $\triangle ABC$ od razpolovišča stranice $AB$.

    
    % \begin{solution}[5cm]

    % \end{solution}











    \newpage
    % 5. vprašanje

    \question
    Daljica $AB$ je določena s krajiščema $A(5,-2,7)$ in $B(1,5,11)$.
    
    \begin{parts}
        \part[3]
        Izračunajte skalarni produkt krajevnega vektorja točke $A$ in vektorja  $\overrightarrow{AB}$.
       
           \begin{solution}[6cm]

    \end{solution}


        \part[3]
        Določite $x$, da bo vektor $\overrightarrow{v}=(2x,7,x+8)$ pravokoten na vektor $\overrightarrow{BA}$.

            \begin{solution}[6cm]

    \end{solution}

        \part[3]
        Določite $y$ in $z$, da bo vektor $\overrightarrow{u}=(z+1,10,y)$ vzporeden krajevnemu vektorju točke $A$.
    \end{parts}

\begin{solution}[4cm]

\end{solution}


\newpage

    % 6. vprašanje



    % \question[4]
    % Dan je trikotnik $\triangle ABC$ s podatki $a=3~cm$, $b=4~cm$ in $c=5~cm$. S kosinusnim izrekom izračunajte velikost kota $\gamma$, ki leži nasproti stranice $c$.
        


    % \begin{solution}[10cm]

    % \end{solution}


\end{questions}



\end{document}